\section{What this review settles}
The background GP and the significance-field GP can be used in the same analysis without circular reasoning merely because both are Gaussian processes. They have different jobs. The first estimates the background inside the moving mask; the second reproduces the correlated fluctuations of the resulting fitted scan. Their agreement with complete Poisson fits tests the latter approximation. It does not independently establish that the spectrum used to generate those fits is a physical background-only source.

The concern about using a source fitted through the observed events is therefore well founded. Such a source can be used for a conditional diagnostic, or within a carefully validated procedure that estimates background parameters from the data. Here, however, the full source is frozen in the toys, and the earlier injection experiment shows substantial absorption when it is rebuilt through an injected peak. A smoother stress response is useful evidence about a particular reference, but is insufficient grounds for adopting that reference as the production null.

\subsection{The quoted 2016 change is more than a trials correction}
The raw observed maximum is still $r=3.4525$ at 90.5 MeV. The reported $Z=2.886$ at 91.5 MeV is the maximum after a mass-dependent transformation, $z_{\rm ref}=(r-a)/s$. Calling these merely two summaries obscured a consequential change of statistical ordering. The reference score is a diagnostic calibration relative to the chosen source; it was not needed just to count search opportunities, and it is not adopted as the primary local result here.

A new audit retains the raw ordering and compares it directly with the saved null maxima. Conditional on the same nominal GP source, the 2016 raw-maximum tail is $0.0678$ in the Gaussian ensemble, with 18 exceedances in 256 complete Poisson scans. The reference-standardized maximum instead gives $0.1049$, with 36/256 direct exceedances. Both are source-conditional probabilities for different tests. The smaller one cannot be selected after inspecting the comparison and then described as an independently specified discovery result.

\subsection{The look-elsewhere correction is not an arbitrary conservatism choice}
A low correlation at one chosen spacing does not divide a continuous search into independent signal regions. A region still contains a maximization, and correlations can return or become negative at greater separations. A \v{S}id\'{a}k calculation becomes a meaningful approximation only after the local and regional test events have been specified and the dependence among those events checked.

The paper's improvement concerns an upcrossing upper bound for the same significance field\cite{AnanievRead}. It does not promise a smaller penalty than counting resolution intervals. The distinction is material for HPS: the fitted response includes sideband subtraction and background retraining, which can produce a more rapidly varying score than a signal template on a known background. The new comparisons below use identical thresholds and covariance matrices wherever possible.

\subsection{What remains fixed and what was newly tested}
All new calculations retain the $\pm2.25\sigma_m$ moving mask. The earlier width scans are summarized as sensitivity tests; they provide no reason to choose a new width now. This review adds a raw-versus-reference maximum audit, a covariance-based test of the proposed ten-percent correlation rule, and a bounded 2021 high-$p_{\rm sum}$ source-transfer diagnostic. It does not repeat the previous full fit scans or claim a new rare-tail calibration.
